\section{When recomputation is mathematically necessary}
\label{sec:C04:invalidation}
The invalidation rule can be stated precisely for a deterministic causal decoder. Let $x_{1:n}$ be its input and let $z_i^{(\ell)}$ be the representation at position $i$ after layer $\ell$. Suppose every layer may attend to every earlier position and its key and value projections are functions of the preceding-layer representation. If position $r$ changes, the original representations at positions $i\geq r$ are not generally reusable. At the input layer, $z_r^{(0)}$ changes. At the first attention layer, every later query may combine a changed value or a changed normalization term from position $r$. At subsequent layers these changed representations produce changed keys and values. Induction across layers establishes a dependency path to every later cached position. This argument establishes possible dependence, not that every numerical value must change for every input.

Exact reuse therefore needs stronger evidence than a small edit distance. A special mask can remove the dependency path; an architecture can partition independent segments; a specific edit may leave a relevant representation unchanged. Each case requires a demonstrated invariant. Ordinary semantic similarity between the original and edited text supplies none of these conditions. The same reasoning applies to changing an adapter or a positional convention. Even if the visible token sequence is identical, the projection functions or their inputs differ. A cache manager should prefer a conservative invalidation rule to silently assuming equality that the model computation does not support.

There is a related distinction between prefix deletion and left truncation during generation. Removing old positions from the attention set while retaining representations computed when those positions were visible does not reproduce a fresh run on the shorter input. It evaluates a hybrid state: later representations already incorporate information from the deleted history, but future queries cannot access the deleted positions directly. Such a policy can be intentional, for example as an approximation that preserves some historical influence. It must be described and evaluated as that policy, rather than as exact execution of the shortened textual context. Training on a freshly shortened sequence would condition on a different computation again.

These observations motivate a three-way equivalence test. First run an uninterrupted reference trajectory. Second resume from a stored cache without editing; this tests serialization and restoration of computational state. Third apply the intended edit and compare with a fresh run under the edited inputs and mask. The second comparison should agree within the documented numerical contract. The third agrees only if the transformation claims exact equivalence; an approximate transformation should instead report its discrepancy and downstream effects. Holding the sampled continuation fixed through teacher forcing helps localize tensor drift before free-running outputs amplify it.

For an agent service, these tests should include a consequential old instruction that is later revoked, not merely irrelevant prose. If removed information survives indirectly in retained states, a textual audit may incorrectly conclude that it has been forgotten. If recomputation eliminates that influence, the cost can be larger than the nominal number of edited tokens suggests. This is a concrete interaction between computational memory, policy conditioning, and deletion semantics. It also explains why the durable event log, the visible context, and the retained cache must be inspected as separate state objects when diagnosing long-session behavior.
